\documentclass[12pt]{article}
\usepackage[utf8]{inputenc}
\usepackage[margin=0.75in]{geometry}
\usepackage{amsmath}\usepackage{amssymb}\usepackage{enumitem}\usepackage{multicol}\usepackage{tcolorbox}
\setlength{\parindent}{0pt}
\begin{document}

\begin{center}{\Large \textbf{Factor a difference of two squares}}\\[2pt]{\small Algebra I \textbar\ CK-12: Factoring Special Products}\end{center}\vspace{0.25cm}

{\large\textbf{Key Idea}}\par\smallskip
A \textbf{difference of squares} $a^2-b^2$ always factors as $(a+b)(a-b)$. Both terms must be perfect squares with a minus between them.\par


{\large\textbf{Key Terms}}\par\smallskip
\begin{itemize}[leftmargin=1.4em]
  \item \textbf{Perfect square} — a number/expression that is something squared
  \item \textbf{Difference of squares} — $a^2-b^2$
\end{itemize}


{\large\textbf{Worked steps}}\par\smallskip
Step 1: Check both terms are perfect squares and separated by a minus.\par
Step 2: Write $(a+b)(a-b)$ where $a,b$ are the square roots.\par
Example: $x^2-9=(x+3)(x-3)$.\par


{\large\textbf{You Try}}\par\smallskip
Factor $x^2-25$.\par

\end{document}